\section{Robustness}\label{app:robust}

\begin{table}[H]
\centering\footnotesize
\caption{Network versus Euclidean indices under alternative specifications (\nCities{} municipalities). ``Mean
gap (\%)'' is the mean of $100\,(I^{N} - I^{E})/I^{E}$; ``Share gap $>$ 0'' is the share of cities in which the
network index is higher. PP: \emph{preta}+\emph{parda}. Population-matched: Euclidean radii shrunk tract by tract
to match the population of 2-km network environments.}\label{tab:robust}
\begin{tabular}{llrrrrr}
\toprule
Specification & Index & Mean gap & Mean gap (\%) & Share gap > 0 & Pearson & Spearman \\
\midrule
2 km, linear & H (five groups) & 0.005 & 29.835 & 1.000 & 0.980 & 0.984 \\
2 km, linear & Dissimilarity (PP) & 0.018 & 12.663 & 0.978 & 0.985 & 0.982 \\
2 km, linear & Isolation (PP) & 0.003 & 0.665 & 0.692 & 0.999 & 0.999 \\
2 km, linear & Entropy (PP) & 0.007 & 28.761 & 1.000 & 0.982 & 0.986 \\
1 km, linear & H (five groups) & 0.007 & 28.578 & 0.989 & 0.984 & 0.988 \\
1 km, linear & Dissimilarity (PP) & 0.017 & 10.965 & 0.978 & 0.989 & 0.984 \\
1 km, linear & Isolation (PP) & 0.004 & 1.152 & 0.802 & 0.999 & 0.999 \\
1 km, linear & Entropy (PP) & 0.008 & 24.604 & 0.978 & 0.987 & 0.988 \\
3 km, linear & H (five groups) & 0.005 & 36.901 & 1.000 & 0.983 & 0.985 \\
3 km, linear & Dissimilarity (PP) & 0.018 & 16.117 & 0.978 & 0.983 & 0.982 \\
3 km, linear & Isolation (PP) & 0.001 & 0.336 & 0.604 & 0.999 & 0.999 \\
3 km, linear & Entropy (PP) & 0.006 & 37.179 & 1.000 & 0.984 & 0.986 \\
2 km, exponential & H (five groups) & 0.005 & 28.481 & 0.989 & 0.983 & 0.986 \\
2 km, exponential & Dissimilarity (PP) & 0.017 & 11.861 & 0.967 & 0.987 & 0.984 \\
2 km, exponential & Isolation (PP) & 0.003 & 0.683 & 0.714 & 0.999 & 0.999 \\
2 km, exponential & Entropy (PP) & 0.006 & 26.950 & 0.989 & 0.985 & 0.988 \\
Population-matched & H (five groups) & -0.000 & -1.577 & 0.209 & 0.999 & 0.999 \\
Population-matched & Dissimilarity (PP) & -0.001 & -0.652 & 0.319 & 0.999 & 0.999 \\
Population-matched & Isolation (PP) & -0.001 & -0.265 & 0.286 & 1.000 & 1.000 \\
Population-matched & Entropy (PP) & -0.000 & -1.382 & 0.253 & 0.999 & 0.999 \\
\bottomrule
\end{tabular}

\end{table}

\begin{table}[H]
\centering\footnotesize
\caption{Street-network topology across the \nCities{} municipalities.}\label{tab:topodesc}
\resizebox{\textwidth}{!}{\begin{tabular}{lrrrrrrrr}
\toprule
 & count & mean & std & min & 25\% & 50\% & 75\% & max \\
\midrule
Intersection density (per km2) & 91.000 & 96.167 & 42.444 & 30.803 & 66.511 & 93.617 & 112.886 & 286.045 \\
Street density (km/km2) & 91.000 & 13.374 & 3.844 & 5.860 & 10.458 & 13.507 & 15.310 & 24.892 \\
Mean street length (m) & 91.000 & 85.027 & 12.286 & 50.293 & 76.284 & 84.081 & 93.440 & 110.996 \\
Streets per node & 91.000 & 2.878 & 0.164 & 2.334 & 2.773 & 2.880 & 2.995 & 3.254 \\
Circuity & 91.000 & 1.045 & 0.016 & 1.018 & 1.033 & 1.043 & 1.054 & 1.112 \\
Dead-end share & 91.000 & 0.149 & 0.052 & 0.054 & 0.116 & 0.145 & 0.178 & 0.360 \\
3-way share & 91.000 & 0.636 & 0.055 & 0.468 & 0.593 & 0.645 & 0.679 & 0.732 \\
4-way share & 91.000 & 0.194 & 0.074 & 0.062 & 0.143 & 0.175 & 0.229 & 0.404 \\
Cyclomatic number & 91.000 & 9371.341 & 11411.163 & 1834.000 & 4142.000 & 6243.000 & 9468.000 & 86122.000 \\
Meshedness & 91.000 & 0.221 & 0.041 & 0.084 & 0.194 & 0.221 & 0.250 & 0.314 \\
Gamma & 91.000 & 0.480 & 0.027 & 0.389 & 0.462 & 0.480 & 0.499 & 0.542 \\
F & 91.000 & 0.552 & 0.152 & 0.283 & 0.420 & 0.542 & 0.680 & 0.896 \\
\bottomrule
\end{tabular}
}
\end{table}

Configurational fragmentation is estimated from random samples of 500 targets and 2,000 origins per city. Doubling
both samples changes $F$ by less than 0.02 in Porto Alegre, Curitiba and Salvador.

\section{Cross-check against the native PySAL network route}\label{app:pandana}

Using the \texttt{pandana} library, the PySAL \texttt{segregation} package can also build network-based
local environments itself \citep{cortes2026segregation}. That route aggregates population at the network node nearest to
each tract and ignores the distance between the tract point and the node. For Porto Alegre, Curitiba and Salvador
it yields network $H$ within 0.001 of our estimates (Table~\ref{tab:pandana}).

\begin{table}[H]
\centering\small
\caption{Network $H$ from our tract-to-tract weights and from \texttt{segregation}'s native \texttt{pandana}
route.}\label{tab:pandana}
\begin{tabular}{lrrr}
\toprule
City & H, Euclidean & H, network (ours) & H, network (pandana) \\
\midrule
Porto Alegre & 0.0647 & 0.0745 & 0.0736 \\
Curitiba & 0.0410 & 0.0476 & 0.0472 \\
Salvador & 0.0310 & 0.0382 & 0.0382 \\
\bottomrule
\end{tabular}

\end{table}

\section{City-level results}\label{app:cities}

{\footnotesize
\begin{longtable}{lrrrr}
\toprule
City & H, Euclidean & H, network & Gap (\%) & F \\
\midrule
\endfirsthead
\toprule
City & H, Euclidean & H, network & Gap (\%) & F \\
\midrule
\endhead
\midrule
\multicolumn{5}{r}{Continued on next page} \\
\midrule
\endfoot
\bottomrule
\endlastfoot
Florianópolis & 0.024 & 0.049 & 108.987 & 0.627 \\
Blumenau & 0.009 & 0.019 & 98.939 & 0.370 \\
Pelotas & 0.016 & 0.027 & 73.877 & 0.668 \\
Juiz de Fora & 0.033 & 0.056 & 67.987 & 0.387 \\
Carapicuíba & 0.002 & 0.003 & 62.676 & 0.341 \\
Ribeirão das Neves & 0.001 & 0.002 & 52.246 & 0.535 \\
Maceió & 0.013 & 0.020 & 52.226 & 0.737 \\
Diadema & 0.005 & 0.008 & 52.102 & 0.541 \\
Serra & 0.008 & 0.012 & 51.846 & 0.794 \\
Ananindeua & 0.003 & 0.004 & 50.797 & 0.604 \\
Ponta Grossa & 0.014 & 0.022 & 50.716 & 0.365 \\
Rio Branco & 0.006 & 0.008 & 50.512 & 0.399 \\
Itaquaquecetuba & 0.002 & 0.003 & 48.926 & 0.643 \\
São Gonçalo & 0.006 & 0.008 & 48.858 & 0.561 \\
Jundiaí & 0.028 & 0.041 & 48.494 & 0.414 \\
Caruaru & 0.003 & 0.004 & 47.485 & 0.308 \\
Jaboatão dos Guararapes & 0.007 & 0.010 & 47.451 & 0.737 \\
Taubaté & 0.010 & 0.015 & 46.754 & 0.283 \\
Paulista & 0.003 & 0.004 & 45.872 & 0.644 \\
Feira de Santana & 0.006 & 0.009 & 44.053 & 0.352 \\
Mauá & 0.011 & 0.016 & 41.736 & 0.388 \\
Vitória & 0.055 & 0.078 & 41.150 & 0.722 \\
Vitória da Conquista & 0.009 & 0.012 & 39.770 & 0.391 \\
São Bernardo do Campo & 0.032 & 0.044 & 39.275 & 0.384 \\
Canoas & 0.012 & 0.016 & 39.116 & 0.504 \\
Belford Roxo & 0.002 & 0.003 & 39.074 & 0.837 \\
Campina Grande & 0.006 & 0.008 & 38.264 & 0.395 \\
São José dos Pinhais & 0.015 & 0.021 & 37.706 & 0.592 \\
Praia Grande & 0.019 & 0.026 & 36.760 & 0.704 \\
Caxias do Sul & 0.024 & 0.033 & 36.372 & 0.357 \\
Macapá & 0.004 & 0.005 & 35.155 & 0.598 \\
Olinda & 0.007 & 0.010 & 31.846 & 0.765 \\
Cariacica & 0.013 & 0.017 & 30.853 & 0.469 \\
Rio de Janeiro & 0.052 & 0.068 & 30.621 & 0.877 \\
Petrolina & 0.008 & 0.011 & 29.791 & 0.414 \\
Recife & 0.031 & 0.040 & 29.134 & 0.719 \\
São Vicente & 0.024 & 0.030 & 29.057 & 0.653 \\
Duque de Caxias & 0.005 & 0.006 & 28.477 & 0.896 \\
Nova Iguaçu & 0.011 & 0.014 & 28.329 & 0.692 \\
Anápolis & 0.013 & 0.016 & 27.118 & 0.454 \\
Belém & 0.012 & 0.015 & 26.333 & 0.585 \\
Contagem & 0.007 & 0.008 & 26.041 & 0.758 \\
Suzano & 0.026 & 0.033 & 25.976 & 0.494 \\
Campos dos Goytacazes & 0.042 & 0.052 & 25.748 & 0.483 \\
Montes Claros & 0.012 & 0.015 & 25.303 & 0.434 \\
Barueri & 0.028 & 0.035 & 25.209 & 0.564 \\
Mogi das Cruzes & 0.030 & 0.037 & 24.837 & 0.378 \\
São Luís & 0.017 & 0.022 & 24.293 & 0.542 \\
Maringá & 0.025 & 0.031 & 23.598 & 0.429 \\
Belo Horizonte & 0.055 & 0.068 & 23.561 & 0.594 \\
Salvador & 0.031 & 0.038 & 23.264 & 0.749 \\
Franca & 0.017 & 0.021 & 23.131 & 0.362 \\
Betim & 0.007 & 0.008 & 22.015 & 0.453 \\
Vila Velha & 0.031 & 0.038 & 21.852 & 0.580 \\
Uberaba & 0.023 & 0.028 & 21.662 & 0.449 \\
Cuiabá & 0.025 & 0.030 & 20.817 & 0.553 \\
Aracaju & 0.017 & 0.020 & 20.423 & 0.397 \\
Sorocaba & 0.023 & 0.027 & 20.049 & 0.337 \\
Niterói & 0.057 & 0.068 & 19.870 & 0.738 \\
Fortaleza & 0.017 & 0.021 & 19.495 & 0.563 \\
João Pessoa & 0.022 & 0.026 & 19.249 & 0.687 \\
Cascavel & 0.035 & 0.042 & 19.029 & 0.695 \\
Uberlândia & 0.028 & 0.034 & 18.436 & 0.439 \\
Goiânia & 0.035 & 0.041 & 17.824 & 0.574 \\
São João de Meriti & 0.004 & 0.005 & 17.682 & 0.785 \\
Porto Velho & 0.010 & 0.012 & 17.604 & 0.469 \\
Londrina & 0.054 & 0.063 & 16.990 & 0.520 \\
Ribeirão Preto & 0.046 & 0.054 & 16.330 & 0.380 \\
Curitiba & 0.041 & 0.048 & 16.257 & 0.435 \\
Piracicaba & 0.051 & 0.059 & 16.229 & 0.412 \\
Joinville & 0.022 & 0.026 & 16.098 & 0.575 \\
Santo André & 0.035 & 0.040 & 16.029 & 0.426 \\
Aparecida de Goiânia & 0.009 & 0.010 & 15.857 & 0.674 \\
Bauru & 0.036 & 0.042 & 15.675 & 0.502 \\
Campinas & 0.067 & 0.077 & 15.247 & 0.502 \\
Porto Alegre & 0.065 & 0.074 & 15.031 & 0.547 \\
Boa Vista & 0.010 & 0.011 & 14.911 & 0.636 \\
Teresina & 0.014 & 0.016 & 14.292 & 0.726 \\
Osasco & 0.018 & 0.021 & 13.971 & 0.430 \\
São José dos Campos & 0.040 & 0.045 & 13.827 & 0.509 \\
Natal & 0.015 & 0.016 & 13.090 & 0.729 \\
Guarulhos & 0.033 & 0.037 & 11.837 & 0.595 \\
São Paulo & 0.072 & 0.079 & 10.791 & 0.502 \\
Campo Grande & 0.035 & 0.039 & 10.096 & 0.569 \\
Palmas & 0.027 & 0.030 & 9.485 & 0.776 \\
Caucaia & 0.026 & 0.028 & 9.060 & 0.695 \\
Santarém & 0.011 & 0.012 & 7.579 & 0.335 \\
São José do Rio Preto & 0.034 & 0.036 & 5.132 & 0.394 \\
Manaus & 0.011 & 0.011 & 4.019 & 0.542 \\
Brasília & 0.043 & 0.044 & 2.716 & 0.844 \\
Santos & 0.041 & 0.041 & 0.751 & 0.809 \\
\end{longtable}

}
